\par\Needspace{9\baselineskip}\pdfbookmark[1]{On the sum of the series}{sixth-section-3}\begin{center}
\textit{THIRD SECTION.}\\[6pt]
\textit{On the sum of our series when the fourth element is set $=1$, with a discussion of certain other}\\
\textit{transcendental functions at the same time.}\\[6pt]
\end{center}
\SixthArticle{15}

Whenever the elements $\alpha,\beta,\gamma$ are all positive quantities, all the coefficients of powers of the fourth element $x$ are positive. If, however, one or another of those elements is negative, all the coefficients will have the same sign at least from some power $x^m$ onwards, provided $m$ is taken greater than the absolute value of the largest negative element. It also follows at once that the sum of the series for
\SourceLeaf{148}
$x=1$ cannot be finite unless the coefficients, at least after a certain term, diminish towards zero, or, to speak in the customary language of analysts, unless the coefficient of the term $x^{\infty}$ is $=0$. We shall show, with full rigour for the benefit of those who favour the rigorous methods of the ancient geometers,

\textit{first}, that the coefficients (unless the series terminates) grow without bound whenever $\alpha+\beta-\gamma-1$ is a positive quantity;

\textit{second}, that the coefficients steadily converge towards a finite limit whenever $\alpha+\beta-\gamma-1=0$;

\textit{third}, that the coefficients diminish towards zero whenever $\alpha+\beta-\gamma-1$ is a negative quantity;

\textit{fourth}, that the sum of our series for $x=1$, despite the convergence in the third case, is infinite whenever $\alpha+\beta-\gamma$ is a positive quantity or $=0$;

\textit{fifth}, that the sum is \textit{finite} whenever $\alpha+\beta-\gamma$ is a negative quantity.

